\scanpage{97}
Donc on s'intéressera \inserted{finalement} au cas où $A_1=k.1+D_1$, $A_2=k.1+D_2$, $D_1,D_2$ contenus dans $M^0=\{x\in M\mid\operatorname{Tr}x=0\}$ et isotropes i.e.
\[
\begin{aligned}
\operatorname{Tr}u&=\operatorname{Tr}v=0\\
\det u&=\det v=0
\end{aligned}
\]

Dans ce cas là on trouve
\[
Q(\unsure{t|b},b,b',c,c')=t^2
\]
Mais dans ce cas là on vérifie aussitôt que pour $x\in M$
\[
\left[\begin{aligned}
\operatorname{Tr}(u\deleted{\unsure{v}}x)=0&\Longleftrightarrow x\text{ normalise }\deleted{\unsure{D|k}}ku\\
\text{i.e.}\qquad [x,u]&=xu-ux=0
\end{aligned}\right.
\]
\deleted{donc} \inserted{i.e.} $x\in B_u$\quad sous-\uuline{alg.} de Borel normalisatrice de $u$,

donc si $t=\operatorname{Tr}(uv)=0$ cela implique que $u,v$ engendrent \deleted{\unsure{un}} une sous-alg. de Borel
\[
k\oplus ku\oplus kv=B_u=B_v
\]
(et \unsure{inversement|nécessairemt}).

Donc
% archive sheet number: 96
